\begin{statement}{C.2.UN.46}{Proved}{The diagonal Hopf cycle criterion and the complete two-component case}
Consider a reduced SNC fibre of an actual smooth-total
degeneration of two-tori, all of whose irreducible components
are unblown diagonal Hopf surfaces of algebraic dimension zero.
Primary and secondary Hopf surfaces are both allowed.
Its components form a cycle: each of the two axis elliptic
curves of each component is a double curve, and there are
no triple points.

For an oriented cycle of specified diagonal Hopf surfaces
$H_1,\ldots,H_\ell$, specify the actual elliptic gluing
isomorphisms and dual-normal isomorphisms along its edges.
Let
\[
 G_i=H_i\setminus(C_i^-\cup C_i^+)
       =\mathbb C^2/\Lambda_i,\qquad \Lambda_i\simeq\mathbb Z^3,
\tag{C.HC.1}
\]
where $\Lambda_i$ is the inverse image of its diagonal
Hopf deck group under coordinatewise exponentiation.
The edge data determine isomorphisms of these integral
lattices and a cycle product $M:\Lambda_1\to\Lambda_1$.
For the finite proper local plumbing constructed below,
the exact criterion is
\[
 \text{$Y_t$ is a two-torus}\quad(0<|t|\ll1)
       \quad\Longleftrightarrow\quad M=I.
\tag{C.HC.2}
\]
Every actual torus degeneration in the stated SNC branch
has $M=I$. The definition of $M$, its independence from
the choices of logarithms, and the actual nearby quotient
are given in the proof.

For two components, the criterion is more explicit.
The two Hopf surfaces are isomorphic. After an actual
diagonal automorphism of one of them, both gluing maps
are inversion on the corresponding axis elliptic curves.
Thus the reduced fibre is precisely two copies of the
same diagonal Hopf surface joined by the two inversion
maps. Every such reduced fibre occurs in a compact
zero-Euler torus filling, by C.2.UN.41 and C.2.UN.43.
This completely determines the reduced supports in
this two-component branch. The constants in the
dual-normal isomorphisms remain explicit local
smoothing parameters in the quotient below.

The same local lattice criterion applies to the
specified operation of identifying the two distinct
axes of one Hopf surface. Its single transition
cannot have $M=I$, so that particular self-gluing
admits no smooth-total torus smoothing.
\end{statement}

\begin{proof}
\par\medskip\noindent\textbf{The cycle and its actual edge data.}
C.2.UN.40 forces both axis curves into the boundary of
every component of the given SNC fibre. An
algebraic-dimension-zero diagonal primary Hopf surface
has exactly these two curves, and its free diagonal
finite quotients have their two distinct images, by
C.2.UN.43. They are disjoint. Consequently there is
no triple point, and at the generic point of either
axis the SNC equation has exactly one other branch.
Every vertex of the connected incidence graph has
degree two, so the graph is a cycle, with a pair
of edges allowed between two vertices.

Orient this cycle. On $H_i$ choose diagonal coordinates
so that the incoming axis $C_i^-$ has coordinate $z_i$
and the outgoing axis $C_i^+$ has coordinate $w_i$.
Thus the outgoing normal coordinate is $z_i$, and
the incoming normal coordinate is $w_i$.
For the edge from $i$ to $i+1$ write its data as
\[
 \phi_i:C_i^+\xrightarrow{\sim}C_{i+1}^-,
 \qquad
 \theta_i:N_{C_i^+/H_i}\otimes
       \phi_i^*N_{C_{i+1}^-/H_{i+1}}
           \xrightarrow{\sim}\mathcal O_{C_i^+}.
\tag{C.HC.3}
\]
These are holomorphic isomorphisms. In an actual
reduced fibre the second comes from the defining
function: locally $f=xy$ times a unit, and its
transverse bilinear term is precisely (C.HC.3).
There is no extra divisor in this isomorphism,
because there are no triple points.

Each axis neighbourhood is an entire holomorphic
normal line bundle. For a primary diagonal Hopf
surface the outgoing example is
\[
 (\mathbb C^*_{w}\times\mathbb C_z)/
       ((w,z)\sim(\beta w,\alpha z)).
\tag{C.HC.4}
\]
For a secondary surface take also the specified
finite diagonal quotient. Its action on the base
elliptic curve is free, since both finite characters
are faithful; the quotient of (C.HC.4) is therefore
an actual line bundle on the quotient elliptic curve.
Removing its zero section gives exactly $G_i$.
Its universal cover is $\mathbb C^2$ in logarithmic
coordinates. The lattice $\Lambda_i$ includes
$2\pi i\mathbb Z^2$ and one logarithm of a contracting
deck transformation, together with the finite
diagonal roots. It is a discrete rank-three lattice,
and it spans $\mathbb C^2$ over $\mathbb C$.

\par\medskip\noindent\textbf{The affine edge maps and the integral product.}
For $t\ne0$, the normal-product equation and
(C.HC.3) define a global biholomorphism
$\Phi_i(t):G_i\to G_{i+1}$: send the base by $\phi_i$
and send a nonzero normal vector to its dual vector
whose product is $t$. Lift it to the logarithmic
universal covers. An elliptic isomorphism lifts to
an affine map, so its base coordinate is
$c_i\log w_i+d_i$, with $c_i\ne0$.
The normal coordinate has the form
$-\log z_i+h_i(\log w_i)+\log t$.
The nonzero holomorphic fibre multiplier has a
holomorphic logarithm $h_i$ on the simply connected
base cover. For every elliptic period $\lambda$,
$h_i(u+\lambda)-h_i(u)$ is constant, since both
normal bundles in (C.HC.4) have constant factors of
automorphy. Hence $h_i'$ descends to a holomorphic
function on a compact elliptic curve and is constant.
We obtain the exact lift
\[
 \binom{\log z_{i+1}}{\log w_{i+1}}
 =
 L_i\binom{\log z_i}{\log w_i}
 +\binom{d_i}{e_i+\log t},
 \qquad
 L_i=\begin{pmatrix}0&c_i\\-1&k_i\end{pmatrix}.
\tag{C.HC.5}
\]
This derivation uses the actual entire line bundles
and their isomorphism; it does not replace an
analytic neighbourhood by a formal normal bundle.

Since $\Phi_i(t)$ is a global biholomorphism,
$L_i\Lambda_i=\Lambda_{i+1}$. Choosing integral bases
gives actual matrices $M_i\in\operatorname{GL}_3(\mathbb Z)$.
Put
\[
 M=M_\ell\cdots M_1,\qquad L=L_\ell\cdots L_1.
\tag{C.HC.6}
\]
Changing the initial basis conjugates $M$.
Changing any logarithm changes only its translation
by a lattice vector. Multiplying $\theta_i$ by a
nonzero constant changes $e_i$, not $L_i$.
Thus the condition $M=I$ is well defined.
It is equivalent to $L=I$, because $\Lambda_1$
spans the complex vector space.

\par\medskip\noindent\textbf{Actual finite plumbing and properness.}
Use (C.HC.4) to choose disjoint tubular neighbourhoods
of the two axes in each Hopf component. They have
flat Hermitian normal metrics. The dual-normal
isomorphism in (C.HC.3) identifies the two flat
metrics up to a positive constant: the ratio of
compatible flat metrics is the exponential of a
harmonic function on the compact elliptic curve.
Adjust that constant so that the product of the
two normal norms is $|t|$ in each neck.

For each edge take the open part of
$N_{C_i^+/H_i}\oplus N_{C_i^+/H_i}^{-1}$
with bounded fibre coordinates $x,y$, and the
actual holomorphic function $xy=t$.
Its two branches are identified with the two
specified axis neighbourhoods using $\phi_i,\theta_i$.
Glue these finitely many necks to the products
of the compact component cores with a disc.
Explicitly choose $0<r_0<R_0$ and $|t|<r_0^2$,
use $\|x\|,\|y\|<R_0$, and attach where one of
these norms is greater than $r_0$.
The inverse map inserts $t/x$ or $t/y$.
The two ends of a neck cannot overlap, because
$\|x\|\|y\|=|t|<r_0^2$; the two axis collars in
a component are disjoint as well.

These are actual biholomorphic gluing maps.
Their graphs are closed relative to the pieces:
at one boundary of an attaching annulus its core
point is absent, and at the other boundary its
neck point is absent. Each equivalence class
meets at most one neck and one core. The finite
gluing is therefore Hausdorff. The intermediate
closed cores and closed smaller neck disc bundles
used in C.2.UN.41 cover the inverse image of every
closed subdisc; each is compact over a compact
elliptic curve or compact Hopf core. This proves
properness. The result is a smooth complex
threefold over a disc, with exactly the prescribed
reduced normal-crossing central surface.

\par\medskip\noindent\textbf{The exact fundamental group, including actual given smoothings.}
The open group $G_i=\mathbb C^2/\Lambda_i$ is
smoothly $T^3\times\mathbb R$.
For example, a linear height annihilating
$\operatorname{span}_{\mathbb R}\Lambda_i$ has
compact three-torus levels; the usual axis norms
are exponentials of this height. A component
core and each annular neck consequently retract
onto a three-torus. Every collar inclusion
induces the indicated isomorphism on its
rank-three fundamental group. Van Kampen,
first along a spanning tree of the cycle and
then at its last edge, gives the actual group
\[
                       \pi_1(Y_t)=\Lambda_1\rtimes_M\mathbb Z.
\tag{C.HC.7}
\]
The copy of $\Lambda_1$ is injected, as in the
displayed semidirect product presentation.

For an arbitrary actual smoothing of the same
reduced SNC data, this calculation still uses
actual smooth maps. Near a double curve $C$,
the divisor line bundles of its two branches
have canonical sections $s_1,s_2$, and $f$
trivializes their tensor product. Choose a
smooth retraction to $C$ and smooth line-bundle
identifications with the pullbacks of their
restrictions, choosing the second identification
so that the product trivialization is preserved.
The map consisting of the retraction and
$(s_1,s_2)$ has invertible derivative along $C$.
After shrinking the neighbourhood it is a
smooth tubular diffeomorphism to the two
normal line bundles, and its actual function
is $xy$. These are the normal bundles and
the product isomorphism (C.HC.3).
On a compact core away from these neighbourhoods
$f$ is a submersion; horizontal transport
identifies its nearby cores. The attaching
maps are isotopic to the normal maps just
specified, by radial contraction of the
tubular identifications. Thus their maps on
$\pi_1$ are exactly $M_i$, proving (C.HC.7)
for the given smoothing as well. No holomorphic
linearization of its full neighbourhood is
asserted or needed.

If the given $Y_t$ is a two-torus, its fundamental
group is abelian. The injected copy of $\Lambda_1$
in (C.HC.7) then forces $M=I$.
This proves the necessary half of (C.HC.2).

\par\medskip\noindent\textbf{Closed one-forms give the sufficient torus comparison.}
Suppose $M=I$, equivalently $L=I$.
The vector of closed holomorphic one-forms
$(d\log z_i,d\log w_i)$ on $G_i$ transforms
by the constant matrix $L_i$ on each overlap.
Choose constant invertible matrices $B_i$
successively so that $B_{i+1}L_i=B_i$.
The identity of the cycle product makes this
choice consistent on its last edge.
These forms therefore glue on each actual
nearby compact fibre, including all its necks,
to two everywhere independent closed holomorphic
one-forms.

Their dual holomorphic vector fields commute,
because the two dual one-forms are closed.
On a compact complex manifold holomorphic
vector fields have entire complex-time flows:
a finite coordinate cover supplies a uniform
local existence interval, which can be iterated.
The two commuting flows give a transitive
holomorphic $\mathbb C^2$-action; every orbit
is open, and the fibre is connected.
Its stabilizer is discrete and cocompact,
since the infinitesimal action is an
isomorphism and the fibre is compact.
The orbit map is therefore the actual
biholomorphism
\[
                  Y_t\simeq\mathbb C^2/\Lambda(t)
\tag{C.HC.8}
\]
for a rank-four real lattice.

One may read this lattice without discarding
the translation data. Compose the affine
maps (C.HC.5) on a chosen logarithm branch:
\[
                    \Phi_\ell(t)\cdots\Phi_1(t)(v)
                          =v+b_0+v_0\log t.
\tag{C.HC.9}
\]
In the normalization $B_1=I$, the period
integrals on the core generators are
$\Lambda_1$ and the cycle generator has
period $b_0+v_0\log t$. Equation (C.HC.7)
lists all generators, so the exact lattice is
\[
                 \Lambda(t)=
                    \Lambda_1+\mathbb Z(b_0+v_0\log t).
\tag{C.HC.10}
\]
Its discreteness and full real rank follow
from the preceding transitive-action proof,
rather than being assumed. A turn of $\log t$
changes the last vector by $2\pi i v_0\in\Lambda_1$,
as follows directly from the globally defined
edge maps. Hence (C.HC.8)--(C.HC.10) descend
to the actual punctured-disc family. This proves
the sufficient half and records every smoothing
translation parameter.

\par\medskip\noindent\textbf{One and two components.}
A single transition in (C.HC.5) cannot be the
identity, because its upper-left entry is zero.
Thus the specified self-identification of the
two distinct axes cannot give a torus smoothing.

For two components multiply the two matrices:
\[
 L_2L_1=
 \begin{pmatrix}
   -c_2&c_2k_1\\
   -k_2&-c_1+k_2k_1
 \end{pmatrix}.
\tag{C.HC.11}
\]
It is the identity precisely when
$c_1=c_2=-1$ and $k_1=k_2=0$.
Consequently $L_1=L_2=-S$, where $S$ exchanges
the two coordinates, and
$\Lambda_2=-S\Lambda_1$.
Exponentiation gives
$\Gamma_2=(S\Gamma_1)^{-1}=S\Gamma_1$,
since these are groups. The actual coordinate
exchange $(z,w)\mapsto(w,z)$ therefore descends
to a biholomorphism $H_1\to H_2$, including
their finite diagonal subgroups.

After this coordinate exchange, the two
elliptic identifications have the form
$z\mapsto a/z$ and $w\mapsto b/w$.
The actual diagonal automorphism
$(Z,W)\mapsto(a^{-1}Z,b^{-1}W)$ of the second
Hopf copy commutes with its complete deck group
and makes both identifications inversion.
It follows that the reduced central surface
is isomorphic to exactly the support in
C.2.UN.41, or its free finite quotient in
C.2.UN.43. Those constructions give the asserted
compact realizations and retain this entire
reduced surface. The constants in $\theta_i$
are recorded by (C.HC.9)--(C.HC.10); the
compact realization assertion here identifies
the reduced support, while the explicit finite
plumbing realizes each specified local choice.
\end{proof}


% C.2.UN.54 from part_C2_hopf_graph_reduction.tex

